\section*{第 \ref{ch:b1:field-lab-methods-1} 章——野外与实验室方法 I}

\begin{solution}{exo:b1:field-lab-methods-1:1}
在 1:10\,000 下，\qty{4.9}{m} 变为 $4.9/10\,000 = \qty{0.49}{mm}$；
在 1:25\,000 时为 $4.9/25\,000 = \qty{0.20}{mm}$，约等于一条铅笔线的宽度。
\end{solution}

\begin{solution}{exo:b1:field-lab-methods-1:2}
在雷克雅未克，$D = \ang{-11.28}$：$52 - 11.28 = \ang{040.7}$。在惠灵顿，
$D = \ang{+23.23}$：$52 + 23.23 = \ang{075.2}$。相同的指南针读数意味着两个方向
相隔 \ang{34.5}。
\end{solution}

\begin{solution}{exo:b1:field-lab-methods-1:3}
$20\to015$。其中 $t = \ang{15}$，$p = \ang{20}$，向下的单位矢量为
\[
  (\cos\ang{20}\sin\ang{15},\ \cos\ang{20}\cos\ang{15},\ -\sin\ang{20}) =
  (0.2432,\ 0.9077,\ -0.3420).
\]
\end{solution}

\begin{solution}{exo:b1:field-lab-methods-1:4}
算术平均为 $(350 + 358 + 6)/3 = \ang{238}$，几乎与
三个读数都相反。水平单位矢量 $(\sin a, \cos a)$ 求和为
$(-0.1736 - 0.0349 + 0.1045,\ 0.9848 + 0.9994 + 0.9945) = (-0.1040, 2.9787)$，
其\omterm{def:b1:spherical-astronomy:horizontal}{方位角}为 $\operatorname{atan2}(-0.1040, 2.9787) = \ang{-2.0}$，也就是
\ang{358.0}：正确的平均值，位于各读数之中。
\end{solution}

\begin{solution}{exo:b1:field-lab-methods-1:5}
均值 $175/5 = \ang{35.0}$。偏差为 $-1, 2, 0, -2, 1$；平方和为 10；
$s = \sqrt{10/4} = \ang{1.58}$。标准误为 $1.58/\sqrt5 = \ang{0.71}$。
\end{solution}

\begin{solution}{exo:b1:field-lab-methods-1:6}
两个法向量分别为 $(\sin\ang{30}, 0, \cos\ang{30}) = (0.5, 0, 0.866)$ 与
$(-0.5, 0, 0.866)$；$\vec S = (0, 0, 1.732)$ 竖直，因此平均为
一个水平面，且 $\bar R = 1.732/2 = 0.866$，每个读数与平均值
相差 \ang{30}。两个读数并非同一面受误差干扰的取样，它们分属
一个\omterm{def:b1:geological-maps:fold}{褶皱}的两翼，是两个不同的总体，二者的平均
值不能描述其中任何一个。每一翼应单独求平均；两者共同确定
\omterm{def:b1:geological-maps:fold}{褶皱轴}，即两个平面的公共直线，此处水平且为南北向。
\end{solution}

\begin{solution}{exo:b1:field-lab-methods-1:7}
$\tan p = \tan\ang{12}\cos(\ang{80} - \ang{120}) = 0.2126\times0.7660 =
0.1628$：$p = \ang{9.2}$，擦痕向 \ang{120} 倾伏。对于
\omterm{def:b1:field-lab-methods-1:trend-plunge}{倾伏向} \ang{170}，$\cos(\ang{-90}) = 0$：该线水平，沿
\omterm{def:b1:geological-maps:attitude}{走向}。
\end{solution}

\begin{solution}{exo:b1:field-lab-methods-1:8}
$0.47/0.20 = 2.35$ 年（自 2026 年 1 月 1 日起算），即 2028 年春季。在
2030 年 1 月 1 日：$-0.47 + 4\times0.20 = \ang{+0.33}$，位于地理北方以东。
\end{solution}

\begin{solution}{exo:b1:field-lab-methods-1:9}
三个法向量位于\omterm{def:b1:spherical-astronomy:horizontal}{方位角}为 \ang{080} 的竖直平面内，与竖直方向的夹角分别为
\ang{10}、\ang{12} 和 \ang{14}；两端的两个关于中间者对称，因此两者之和
沿该方向：平均为
$080/12$，其中 $\theta = \ang{2}, 0, \ang{2}$。于是 $\bar R = (1 +
2\cos\ang{2})/3 = (1 + 1.998782)/3 = 0.999594$，且 $\sqrt{2(1 - \bar R)} =
\sqrt{0.000812} = \qty{0.02850}{rad} = \ang{1.633}$；精确地，
$\theta_{\text{方均根}} = \sqrt{(4 + 0 + 4)/3} = \ang{1.633}$。
\end{solution}

\begin{solution}{exo:b1:field-lab-methods-1:10}
截面为半径 $\sqrt{r^2 - h^2}$ 的圆。其平均值为
$\frac1r\int_0^r\sqrt{r^2 - h^2}\,\mathrm dh = \frac1r\cdot\frac{\pi r^2}{4}
= \frac{\pi r}{4}$，因为该积分即四分之一圆盘的面积。
颗粒宽 \qty{0.50}{mm}，平均视宽为 $0.50\times\pi/4 = \qty{0.39}{mm}$，
\omterm{def:b1:field-lab-methods-1:thin-section}{薄片}会低估颗粒的大小。
\end{solution}

\begin{solution}{exo:b1:field-lab-methods-1:11}
绕竖直方向转过角 $D$ 的旋转，
$(x, y, z) \mapsto (x\cos D + y\sin D,\ -x\sin D + y\cos D,\ z)$，将
$(\sin\delta\sin\alpha, \sin\delta\cos\alpha, \cos\delta)$ 映为
$(\sin\delta\sin(\alpha + D), \sin\delta\cos(\alpha + D), \cos\delta)$：它
使每个\omterm{def:b1:geological-maps:attitude}{倾向}都增加 $D$。由于线性，该变换将合矢量
$\vec S$ 映为转动后各法向量的合矢量；它保持长度和
竖直分量，因此平均\omterm{def:b1:geological-maps:attitude}{倾角} $\arccos(S_U/\|\vec S\|)$ 不变，
并将 $\vec S$ 的水平分量转过 $D$，因此平均\omterm{def:b1:geological-maps:attitude}{倾向}
增加 $D$。对每个读数，或只对它们的
平均值作磁偏角修正，所得结果相同。
\end{solution}

\begin{solution}{exo:b1:field-lab-methods-1:12}
沿最大倾斜线方向的离散，即\omterm{def:b1:geological-maps:attitude}{倾角}的离散：
标准误为 $2/\sqrt N \le 0.5$，当 $N \ge 16$ 时。与之垂直的方向上，法线的
转动 \ang{2} 使\omterm{def:b1:geological-maps:attitude}{倾向}改变 $2/\sin\ang{12} = \ang{9.62}$；
平均\omterm{def:b1:geological-maps:attitude}{倾向}的标准误为 $9.62/\sqrt N \le 1$，当
$N \ge 92.5$ 时，即需 93 个读数。平缓的\omterm{def:b1:geological-maps:attitude}{倾角}很快便可确定；其
\omterm{def:b1:geological-maps:attitude}{倾向}则不然。
\end{solution}

\begin{solution}{pb:b1:field-lab-methods-1:1}
\textbf{1.} \omterm{def:b1:geological-maps:attitude}{倾向}：$967/12 = \ang{80.6}$；\omterm{def:b1:geological-maps:attitude}{倾角}：$145/12 =
\ang{12.1}$。

\textbf{2.} \ang{355} 与 \ang{005} 的平均值为 \ang{180}，即相反的
方向。这里所有\omterm{def:b1:geological-maps:attitude}{倾向}都位于 \ang{076} 与 \ang{085} 之间，
远离 \ang{000} 处的跳跃，因此算术平均可用。

\textbf{3.} $\|\vec n\|^2 =\allowbreak \sin^2\delta(\sin^2\alpha + \cos^2\alpha) +
\cos^2\delta = 1$。向下的倾斜线方向为
$(\cos\delta\sin\alpha,\allowbreak \cos\delta\cos\alpha,\allowbreak -\sin\delta)$，其
与 $\vec n$ 的点积为 $\sin\delta\cos\delta - \cos\delta\sin\delta = 0$；\omterm{def:b1:geological-maps:attitude}{走向}
线水平，且垂直于\omterm{def:b1:geological-maps:attitude}{倾向}，
其方向为 $(\cos\alpha, -\sin\alpha, 0)$，它与 $\vec n$ 的点积为
$\sin\delta\sin\alpha\cos\alpha - \sin\delta\cos\alpha\sin\alpha = 0$。

\textbf{4.} 读数 11 与读数 2 相同：$(0.2222, 0.0352, 0.9744)$。
读数 12：
\[
  (\sin\ang{11}\sin\ang{82},\ \sin\ang{11}\cos\ang{82},\ \cos\ang{11}) =
  (0.1890,\ 0.0266,\ 0.9816).
\]

\textbf{5.} $\vec S = (2.4747,\ 0.4082,\ 11.7294)$。

\textbf{6.} $\|\vec S\| = \sqrt{6.1241 + 0.1666 + 137.5788} = 11.99457$，且
$\bar R = 11.99457/12 = 0.999547$。

\textbf{7.} $\bar\alpha = \operatorname{atan2}(2.4747, 0.4082) = \ang{80.6}$
且 $\bar\delta = \arccos(11.7294/11.99457) = \arccos 0.977893 = \ang{12.1}$：
矢量平均为 $080.6/12.1$，与此处的算术平均相同。

\textbf{8.} \omterm{def:b1:geological-maps:attitude}{倾向}：$s = \ang{2.94}$；\omterm{def:b1:geological-maps:attitude}{倾角}：$s = \ang{1.68}$。

\textbf{9.} $\theta_{\text{方均根}} \approx \sqrt{2\times0.000453} =
\qty{0.0301}{rad} = \ang{1.72}$。

\textbf{10.} $\sqrt{1.68^2 + (\sin\ang{12.1}\times2.94)^2} = \sqrt{2.82 +
0.38} = \ang{1.79}$，接近 \ang{1.72}：读数的表现与误差模型所述相符，
其偏差在按整度取值与
样本量较小的影响范围之内。

\textbf{11.} \omterm{def:b1:geological-maps:attitude}{倾角}的贡献：它占 3.20（以平方度计）中的 2.82。
\omterm{def:b1:geological-maps:attitude}{倾向}只通过 $\sin\delta = 0.21$ 起作用，因为在较缓的
岩层上，\omterm{def:b1:geological-maps:attitude}{倾向}的变化使平面的转动很小。

\textbf{12.} 平均\omterm{def:b1:geological-maps:attitude}{倾角}：$1.68/\sqrt{12} = \ang{0.48}$；平均\omterm{def:b1:geological-maps:attitude}{倾向}：
$2.94/\sqrt{12} = \ang{0.85}$。

\textbf{13.} 否：\omterm{def:b1:geological-maps:attitude}{倾角}为 \ang{12} 的岩层上，\omterm{def:b1:geological-maps:attitude}{倾向}的 \ang{2.9} 仅相当于
$2.9\sin\ang{12.1} = \ang{0.6}$ 的平面转动，小于
\ang{1.7} 的\omterm{def:b1:geological-maps:attitude}{倾角}读数。岩层面的任何轻微起伏
都会使其\omterm{def:b1:geological-maps:attitude}{倾向}大幅分散。

\textbf{14.} $D = \ang{-2.5}$：$83.1 - 2.5 = \ang{080.6}$；她的平均值为
$080.6/12.1$，与第一位地质学家的相同。

\textbf{15.} \omterm{def:b1:geological-maps:attitude}{倾角}用测斜仪读取，以重力给出的
竖直方向为基准；磁针在此不起作用。

\textbf{16.} 是：绕竖直方向的转动会使每个读数及
平均值都转过相同的角度（见\cref{exo:b1:field-lab-methods-1:11}）；答案不变。

\textbf{17.} 每个\omterm{def:b1:geological-maps:attitude}{倾向}都有 $10\times0.2 = \ang{2}$ 的误差，
超过 \ang{0.85} 标准误差的两倍：这是一种系统误差，
平均无法将其检出。

\textbf{18.} $80.6 - 90 = \ang{350.6}$；\omterm{def:b1:geological-maps:attitude}{产状}为\omterm{def:b1:geological-maps:attitude}{走向} \ang{351}，
\omterm{def:b1:geological-maps:attitude}{倾向}东，\omterm{def:b1:geological-maps:attitude}{倾角} \ang{12}。

\textbf{19.} \omterm{def:b1:geological-maps:attitude}{倾向} $080.6$ 与 $080.6$：在此精度下无差别；
\omterm{def:b1:geological-maps:attitude}{倾角} $12.1$ 与 $12.2$：相差 \ang{0.1}。

\textbf{20.} 法向量 $(\sin\ang{12.2}\sin\ang{80.6},
\sin\ang{12.2}\cos\ang{80.6}, \cos\ang{12.2})$ 与 $\vec S/\|\vec S\|$ 的夹角为
\ang{0.13}。

\textbf{21.} $1.72/\sqrt{12} = \ang{0.50}$；差值 \ang{0.13} 为其
四分之一，差异不显著。指南针与钻孔结果一致。

\textbf{22.} 数量为四倍：共 48 个读数。

\textbf{23.} 例如：磁偏角设置错误或已过时；指南针
受附近铁质物体（车辆、工具、磁性岩石）影响而偏转；测斜仪
校准不当；仅在露出最好的面上取读数，这些面
可能无法代表该翼。平均可以减小随机误差，而不能减小系统
误差。

\textbf{24.} 东翼的平均\omterm{def:b1:geological-maps:attitude}{产状}为 \textbf{081/12}：\omterm{def:b1:geological-maps:attitude}{倾向}
\ang{081}，\omterm{def:b1:geological-maps:attitude}{倾角} \ang{12}。
\end{solution}
